For $N\ge2$, let $p_i$ be the probability of reaching $N$ first when starting at $i$. Then $p_0=0,p_N=1$, and conditioning on the first move gives
\[
p_i=\dfrac{p_{i-1}+p_{i+1}}{2},\qquad1\le i<N.
\]
Therefore the consecutive differences $p_i-p_{i-1}$ are constant. The boundary values force $p_i=\dfrac{i}{N}$. Absorption occurs almost surely: from any interior state, $N$ consecutive heads suffice to reach the upper boundary, so the probability of avoiding absorption through $k$ blocks of $N$ tosses is at most $(1-2^{-N})^k$, tending to zero.

Thus (a) $\boxed{P(3)=\dfrac{1}{3}}$ and (b) $\boxed{P(N)=\dfrac{1}{N}}$. The latter also gives the immediate-stop value $P(1)=1$.
